% Part: first-order-logic
% Chapter: introduction
% Section: models-theories

\documentclass[../../../include/open-logic-section]{subfiles}

\begin{document}

\olfileid[pt]{fol}{int}{mod}

\olsection{Modelos e Teorias}

Uma vez que definimos a sintaxe e a semântica da lógica de primeira
ordem, podemos trabalhar investigando as propriedades de
!!{structure}s e as noções semânticas. Também podemos definir sistemas
de !!{derivation} e investigá-los.  Para um conjunto de
!!{sentence}s, podemos perguntar: quais !!{structure}s tornam todas as
!!{sentence}s desse conjunto verdadeiras?  Dado um conjunto de
!!{sentence}s~$\Gamma$, uma !!{structure}~$\Struct{M}$ que as satisfaz
é chamada de \emph{modelo de~$\Gamma$}.  Podemos começar
de~$\Gamma$ e tentar encontrar seus modelos---como eles são? Quão
grandes ou pequenos precisam ser? Mas também podemos começar com uma
única !!{structure} ou coleção de !!{structure}s e perguntar: quais
!!{sentence}s são verdadeiras nelas?  Há !!{sentence}s que
\emph{caracterizam} essas !!{structure}s no sentido de que elas, e
somente elas, são verdadeiras nelas? Esses tipos de questões são o
domínio da \emph{teoria dos modelos}.  Elas também fundamentam o
\emph{método axiomático}: descrever uma coleção de !!{structure}s por
um conjunto de !!{sentence}s, os axiomas de uma teoria. Isso é possível
pela observação de que exatamente aquelas !!{sentence}s que são consequência
na lógica de primeira ordem dos axiomas são verdadeiras em todos os
modelos dos axiomas.

Como exemplo muito simples, considere pré-ordens. Uma pré-ordem é uma
relação~$R$ sobre algum conjunto~$A$ que é reflexiva e transitiva.
Um conjunto~$A$ com uma relação binária $R \subseteq A \times A$ sobre
ele é exatamente o que precisaríamos para dar uma !!{structure} para uma
linguagem de primeira ordem com um único símbolo de predicado de dois
lugares~$\Obj P$: definiríamos $\Domain{M} = A$ e $\Assign{\Obj P}{M} =
R$.  Como $R$~é uma pré-ordem, é reflexiva e transitiva, e podemos
encontrar um conjunto~$\Gamma$ de !!{sentence}s da lógica de primeira
ordem que diz isso:
\begin{align*}
  & \lforall[\Obj v_0][\Atom{\Obj P}{\Obj v_0,\Obj v_0}]\\
  & \lforall[\Obj v_0][\lforall[\Obj v_1][\lforall[\Obj v_2][((\Atom{\Obj P}{\Obj v_0,\Obj v_1} \land \Atom{\Obj P}{\Obj v_1,\Obj v_2}) \lif \Atom{\Obj P}{\Obj v_0,\Obj v_2})]]]
\end{align*}
Estas !!{sentence}s são apenas as simbolizações de ``para todo~$x$,
$Rxx$'' ($R$ é reflexiva) e ``sempre que $Rxy$ e $Ryz$, então também
$Rxz$'' ($R$ é transitiva). Vemos que uma !!{structure}~$\Struct{M}$
é um modelo dessas duas !!{sentence}s~$\Gamma$ se, e somente se, $R$
(isto é, $\Assign{\Obj P}{M}$), é uma pré-ordem sobre~$A$ (isto é,
$\Domain{M}$). Em outras palavras, os modelos de $\Gamma$ são
exatamente as pré-ordens. Qualquer propriedade de todas as pré-ordens
que possa ser expressa na linguagem de primeira ordem apenas com~$\Obj P$
como !!{predicate} (como reflexividade e transitividade acima), é
consequência das duas !!{sentence}s em~$\Gamma$ e vice-versa.  Assim,
qualquer coisa que possamos provar sobre modelos de~$\Gamma$ provamos
sobre todas as pré-ordens.

Para qualquer teoria particular e classe de modelos (como $\Gamma$ e
todas as pré-ordens), haverá questões interessantes sobre o que pode ser
expresso na linguagem de primeira ordem correspondente, e o que não pode
ser expresso. Há algumas propriedades de !!{structure}s que são
interessantes para todas as linguagens e classes de modelos, a saber,
aquelas que dizem respeito ao tamanho do !!{domain}.  Pode-se sempre
expressar, por exemplo, que o !!{domain} contém exatamente $n$~!!{element}s,
para qualquer~$n \in \PosInt$.  Pode-se também expressar, usando um
conjunto de infinitas !!{sentence}s, que o !!{domain} é
infinito.  Mas não se pode expressar que o domínio é finito, ou que o
domínio é !!{nonenumerable}. Estes resultados sobre as limitações das
linguagens de primeira ordem são consequências dos teoremas da
compacidade e de L\"owenheim--Skolem.

\end{document}
